\chapter{Wymagania systemowe}

\section{Wymagania funkcjonalne aplikacji}

Aplikacja powinna umożliwiać podstawową obsługę danych hotelowych. System nie jest pełnym programem komercyjnym, tylko prostą aplikacją do zarządzania najważniejszymi elementami hotelu.

\begin{table}[H]
\centering
\caption{Najważniejsze wymagania funkcjonalne aplikacji}
\begin{tabular}{|p{3cm}|p{10cm}|}
\hline
\textbf{Moduł} & \textbf{Wymaganie} \\
\hline
Goście & Dodawanie nowych gości oraz wyświetlanie ich w tabeli. \\
\hline
Pokoje & Przeglądanie pokoi, numerów, statusów oraz typów pokoi. \\
\hline
Rezerwacje & Tworzenie rezerwacji dla wybranego gościa i pokoju. \\
\hline
Dostępność & Sprawdzanie, które pokoje są dostępne w wybranym terminie. \\
\hline
Płatności & Dodawanie płatności do wybranej rezerwacji. \\
\hline
Usługi & Dodawanie i edytowanie usług dodatkowych. \\
\hline
Panel główny & Wyświetlanie statystyk, liczby gości, rezerwacji i płatności. \\
\hline
\end{tabular}
\end{table}

Do wymagań funkcjonalnych należą także walidacja podstawowych pól, odświeżanie danych po wykonaniu operacji oraz pokazywanie komunikatów o błędach w dolnej części okna.

\section{Wymagania niefunkcjonalne}

Aplikacja powinna działać lokalnie na komputerze z systemem Windows. Interfejs powinien być czytelny i spójny, aby użytkownik mógł łatwo przechodzić między modułami. Kod powinien być podzielony na warstwy i pliki według odpowiedzialności.

Do wymagań niefunkcjonalnych należą:

\begin{itemize}
  \item lokalne działanie programu na komputerze z systemem Windows,
  \item zapis danych w lokalnej bazie SQLite,
  \item czytelny układ głównego okna,
  \item podział projektu na modele, kontekst bazy, repozytoria, serwisy i ViewModel,
  \item podstawowa obsługa wyjątków,
  \item walidacja danych wpisywanych w formularzach,
  \item użycie Entity Framework Core do dostępu do danych,
  \item możliwość łatwego uruchomienia projektu w Visual Studio 2022.
\end{itemize}

\section{Założenia dotyczące działania systemu}

Założono, że program działa lokalnie i korzysta z jednej bazy danych zapisanej w pliku \texttt{hotel.db}. Użytkownik nie musi zakładać konta i nie musi się logować. Dane testowe są dodawane automatycznie przy pierwszym uruchomieniu aplikacji.

System służy do podstawowego zarządzania hotelem. Nie zawiera modułu logowania, faktur, rezerwacji online ani integracji z zewnętrznymi serwisami. Program można jednak rozbudować w przyszłości o takie elementy.
